\section{A Complete FRI Folding Example over \texorpdfstring{$\F_{17}$}{F17}}
\label{app:fri-example}

This appendix gives a complete arithmetic transcript for binary multiplicative-domain FRI
protocol of Section~\ref{sec:fri}. The field and domain are far too small for cryptographic
security. Their purpose is to make every evaluation, fold, terminal degree check, and Merkle
opening location explicit.

\subsection{Parameters and evaluation domains}

Work over
\[
    \F=\F_{17}.
\]
The element
\[
    \omega=3
\]
has order $16$ in $\F_{17}^{\times}$. Take the initial domain
\[
    D_0=\angles{\omega}
       =(1,3,9,10,13,5,15,11,16,14,8,7,4,12,2,6).
\]
The order shown here is important: it is the leaf order used by the first Merkle tree.
Squaring gives
\[
    D_1=\angles{\omega^2}
       =\angles{9}
       =(1,9,13,15,16,8,4,2),
\]
and squaring once more gives
\[
    D_2=\angles{\omega^4}
       =\angles{13}
       =(1,13,16,4).
\]
The sizes are
\[
    |D_0|=16,
    \qquad
    |D_1|=8,
    \qquad
    |D_2|=4.
\]

We prove proximity to the Reed--Solomon code
\[
    \RS[\F_{17},D_0,4].
\]
Thus the claimed degree bound is
\[
    d_0=4
\]
and the code rate is
\[
    \rho=\frac{4}{16}=\frac14.
\]
Each fold halves both the domain size and the degree bound:
\[
\begin{array}{c|c|c|c}
\text{round }i & |D_i| & d_i & d_i/|D_i|\\
\hline
0 & 16 & 4 & 1/4\\
1 & 8 & 2 & 1/4\\
2 & 4 & 1 & 1/4
\end{array}
\]
After two folds, a degree-$<1$ polynomial must be constant.

\subsection{The initial Reed--Solomon word}

Let the initial polynomial be
\[
    p_0(X)=3+5X+2X^2+7X^3.
\]
It has degree $3<4$, so its evaluation word is a valid codeword. Evaluating in the domain
order above gives
\[
\begin{aligned}
\mathbf f^{(0)}
&=\bigl(p_0(x)\bigr)_{x\in D_0}\\
&=(0,4,9,11,9,1,13,12,10,4,15,4,10,3,9,2).
\end{aligned}
\]
For example,
\[
\begin{aligned}
p_0(10)
&=3+5\cdot10+2\cdot10^2+7\cdot10^3\\
&=7253\\
&\equiv 11\pmod{17}.
\end{aligned}
\]
The points halfway around the subgroup differ by a sign because
\[
    \omega^8=16=-1\pmod{17}.
\]
Thus the two preimages under squaring are paired by indices $j$ and $j+8$:
\[
\begin{array}{c|c|c}
x & -x & x^2\\
\hline
1 & 16 & 1\\
3 & 14 & 9\\
9 & 8 & 13\\
10 & 7 & 15\\
13 & 4 & 16\\
5 & 12 & 8\\
15 & 2 & 4\\
11 & 6 & 2
\end{array}
\]

\subsection{First folding round}

Split $p_0$ into even and odd coefficient polynomials:
\[
\begin{aligned}
p_0(X)
&=(3+2X^2)+X(5+7X^2)\\
&=p_{0,\mathrm{even}}(X^2)
  +X p_{0,\mathrm{odd}}(X^2),
\end{aligned}
\]
where
\[
    p_{0,\mathrm{even}}(Y)=3+2Y,
    \qquad
    p_{0,\mathrm{odd}}(Y)=5+7Y.
\]
Let the first verifier challenge be
\[
    \alpha_0=4.
\]
The folded polynomial is
\[
\begin{aligned}
p_1(Y)
&=p_{0,\mathrm{even}}(Y)+\alpha_0p_{0,\mathrm{odd}}(Y)\\
&=(3+2Y)+4(5+7Y)\\
&=23+30Y\\
&\equiv 6+13Y\pmod{17}.
\end{aligned}
\]
As required,
\[
    \deg p_1=1<2=\frac{d_0}{2}.
\]

We now reproduce this fold using only evaluation values. Since
\[
    2^{-1}\equiv9\pmod{17},
\]
define for each pair $\{x,-x\}$
\[
    E_0(x^2)=\frac{f^{(0)}(x)+f^{(0)}(-x)}2,
\]
\[
    O_0(x^2)=\frac{f^{(0)}(x)-f^{(0)}(-x)}{2x}.
\]
The complete first-round table is:
\[
\begin{array}{c|c|c|c|c|c|c|c}
x & -x & y=x^2
& f^{(0)}(x) & f^{(0)}(-x)
& E_0(y) & O_0(y) & E_0(y)+4O_0(y)\\
\hline
1  &16&1 &0 &10&5 &12&2\\
3  &14&9 &4 &4 &4 &0 &4\\
9  &8 &13&9 &15&12&11&5\\
10 &7 &15&11&4 &16&8 &14\\
13 &4 &16&9 &10&1 &15&10\\
5  &12&8 &1 &3 &2 &10&8\\
15 &2 &4 &13&9 &11&16&7\\
11 &6 &2 &12&2 &7 &2 &15
\end{array}
\]
These folded values are exactly the evaluations of $p_1(Y)=6+13Y$ on $D_1$:
\[
\begin{aligned}
\mathbf f^{(1)}
&=\bigl(p_1(y)\bigr)_{y\in D_1}\\
&=(2,4,5,14,10,8,7,15).
\end{aligned}
\]

As one nontrivial row check, take $x=9$, $-x=8$, and $y=x^2=13$. We have
\[
    f^{(0)}(9)=9,
    \qquad
    f^{(0)}(8)=15.
\]
The even component is
\[
    E_0(13)=\frac{9+15}{2}=24\cdot9\equiv12,
\]
and, because
\[
    (2\cdot9)^{-1}=1^{-1}=1\pmod{17},
\]
the odd component is
\[
    O_0(13)=\frac{9-15}{2\cdot9}\equiv -6\equiv11.
\]
Therefore
\[
    E_0(13)+4O_0(13)=12+4\cdot11=56\equiv5,
\]
which agrees with
\[
    p_1(13)=6+13\cdot13=175\equiv5.
\]

\subsection{Second folding round}

The degree-$1$ polynomial $p_1$ already has the decomposition
\[
    p_1(Y)=6+13Y
          =p_{1,\mathrm{even}}(Y^2)
           +Yp_{1,\mathrm{odd}}(Y^2),
\]
where
\[
    p_{1,\mathrm{even}}(Z)=6,
    \qquad
    p_{1,\mathrm{odd}}(Z)=13.
\]
Let the second verifier challenge be
\[
    \alpha_1=6.
\]
Then
\[
\begin{aligned}
p_2(Z)
&=p_{1,\mathrm{even}}(Z)+\alpha_1p_{1,\mathrm{odd}}(Z)\\
&=6+6\cdot13\\
&=84\\
&\equiv16\pmod{17}.
\end{aligned}
\]
Thus $p_2$ is the constant polynomial $16$, as required by the terminal degree bound
$d_2=1$.

The complete evaluation-level fold is:
\[
\begin{array}{c|c|c|c|c|c|c|c}
y & -y & z=y^2
& f^{(1)}(y) & f^{(1)}(-y)
& E_1(z) & O_1(z) & E_1(z)+6O_1(z)\\
\hline
1  &16&1 &2 &10&6&13&16\\
9  &8 &13&4 &8 &6&13&16\\
13 &4 &16&5 &7 &6&13&16\\
15 &2 &4 &14&15&6&13&16
\end{array}
\]
Therefore
\[
    \mathbf f^{(2)}
    =\bigl(p_2(z)\bigr)_{z\in D_2}
    =(16,16,16,16).
\]
In the terminal step the prover need not commit to four identical values. It can send the
single coefficient
\[
    c_2=16,
\]
and the verifier treats the terminal oracle as the constant function $f^{(2)}(z)=c_2$.

\subsection{A concrete commit transcript}

We now place the arithmetic tables inside the cryptographic transcript. Let
$\mathsf{MerkleRoot}$ denote a binary Merkle-tree root, with the leaf index included in the
hashed leaf encoding.

The prover first commits to the initial evaluation vector:
\[
    R_0=\mathsf{MerkleRoot}\bigl(\mathbf f^{(0)}\bigr).
\]
The verifier sends $\alpha_0=4$. In a non-interactive proof this value would be derived as
\[
    \alpha_0
    =\mathsf{HashToField}
      \bigl(\mathsf{fri\mbox{-}alpha\mbox{-}0}\,\|\,R_0\bigr).
\]
For this toy transcript, we stipulate that the resulting field element is $4$.

The prover computes and commits to the first folded vector:
\[
    R_1=\mathsf{MerkleRoot}\bigl(\mathbf f^{(1)}\bigr).
\]
The second challenge is
\[
    \alpha_1
    =\mathsf{HashToField}
      \bigl(\mathsf{fri\mbox{-}alpha\mbox{-}1}\,\|\,R_0\,\|\,R_1\bigr)
    =6.
\]
The prover computes the terminal polynomial and sends
\[
    c_2=16.
\]
Only now is the query index derived:
\[
\begin{aligned}
j
&=\mathsf{HashToIndex}
  \bigl(\mathsf{fri\mbox{-}query}\,\|\,R_0\,\|\,R_1\,\|\,c_2\bigr)\\
&=0.
\end{aligned}
\]
The equalities to $4$, $6$, and $0$ are representative Fiat--Shamir outputs chosen to make
the arithmetic transcript concrete. They are not claims about a particular deployed hash
function.

\subsection{One complete query path}

The query index $j=0$ selects
\[
    x_0=\omega^0=1\in D_0.
\]
The path of squared points is
\[
    x_0=1,
    \qquad
    x_1=x_0^2=1\in D_1,
    \qquad
    x_2=x_1^2=1\in D_2.
\]

\subsubsection*{Round 0 check}

The verifier opens
\[
    f^{(0)}(1)=0,
    \qquad
    f^{(0)}(-1)=f^{(0)}(16)=10,
\]
under $R_0$, and
\[
    f^{(1)}(1)=2
\]
under $R_1$. The locally reconstructed even value is
\[
    \frac{0+10}{2}=10\cdot9=90\equiv5,
\]
and the odd value is
\[
    \frac{0-10}{2\cdot1}
    =7\cdot9
    =63
    \equiv12.
\]
The FRI equation gives
\[
\begin{aligned}
\frac{f^{(0)}(1)+f^{(0)}(16)}2
+\alpha_0
 \frac{f^{(0)}(1)-f^{(0)}(16)}{2}
&=5+4\cdot12\\
&=53\\
&\equiv2\\
&=f^{(1)}(1).
\end{aligned}
\]
The first round accepts.

\subsubsection*{Round 1 check}

The already opened value $f^{(1)}(1)=2$ is reused. The verifier additionally opens
\[
    f^{(1)}(-1)=f^{(1)}(16)=10
\]
under $R_1$. The even and odd values are
\[
    \frac{2+10}{2}=12\cdot9=108\equiv6,
\]
\[
    \frac{2-10}{2\cdot1}=9\cdot9=81\equiv13.
\]
Using $\alpha_1=6$,
\[
\begin{aligned}
\frac{f^{(1)}(1)+f^{(1)}(16)}2
+\alpha_1
 \frac{f^{(1)}(1)-f^{(1)}(16)}2
&=6+6\cdot13\\
&=84\\
&\equiv16\\
&=c_2.
\end{aligned}
\]
The second round accepts, and $c_2$ describes a degree-$<1$ terminal polynomial. This
completes the arithmetic verification path.

\subsection{The Merkle authentication paths}

To make the oracle realization explicit, let $L_k^{(i)}$ denote leaf $k$ of the tree for
$\mathbf f^{(i)}$, and let $N_{a:b}^{(i)}$ denote the internal node covering the contiguous
leaf interval $a,\ldots,b$.

In the first tree, the round-$0$ query opens indices
\[
    0\quad\text{and}\quad8,
\]
corresponding to the points $1$ and $16=-1$. A separate authentication path for index $0$
would contain
\[
    L_1^{(0)},\quad
    N_{2:3}^{(0)},\quad
    N_{4:7}^{(0)},\quad
    N_{8:15}^{(0)}.
\]
A separate path for index $8$ would contain
\[
    L_9^{(0)},\quad
    N_{10:11}^{(0)},\quad
    N_{12:15}^{(0)},\quad
    N_{0:7}^{(0)}.
\]
When the two openings are batched, $N_{0:7}^{(0)}$ and $N_{8:15}^{(0)}$ are reconstructed
from the opened branches. The minimal sibling-node set is therefore
\[
\set{
L_1^{(0)},
N_{2:3}^{(0)},
N_{4:7}^{(0)},
L_9^{(0)},
N_{10:11}^{(0)},
N_{12:15}^{(0)}
}.
\]
Together with the opened leaf values $0$ and $10$, these nodes reconstruct $R_0$.

In the second tree, the full path needs indices
\[
    0\quad\text{and}\quad4,
\]
corresponding to $1$ and $16=-1$ in the $D_1$ ordering. The batched sibling-node set is
\[
\set{
L_1^{(1)},
N_{2:3}^{(1)},
L_5^{(1)},
N_{6:7}^{(1)}
}.
\]
Together with the opened values $2$ and $10$, these nodes reconstruct $R_1$.

The verifier rejects before doing field arithmetic if either reconstructed root differs from
the committed root. Merkle verification therefore binds the numerical inputs used in both
folding equations to the oracles fixed during the commit phase.

\subsection{How a corrupted value is detected}

Suppose, for illustration, that the first entry were changed from
\[
    f^{(0)}(1)=0
\]
to
\[
    \widetilde f^{(0)}(1)=1,
\]
while the queried values $f^{(0)}(16)=10$ and $f^{(1)}(1)=2$ remained unchanged. Under the
same challenge $\alpha_0=4$, the verifier would reconstruct
\[
    \widetilde E_0(1)
    =\frac{1+10}{2}
    =11\cdot9
    =99
    \equiv14,
\]
\[
    \widetilde O_0(1)
    =\frac{1-10}{2}
    =8\cdot9
    =72
    \equiv4.
\]
The expected folded value would become
\[
    14+4\cdot4=30\equiv13,
\]
but the committed next-round value is $2$. Hence
\[
    13\neq2,
\]
and this query path rejects.

This calculation illustrates local consistency detection. It is not by itself a soundness
proof: a malicious prover is free to alter later oracles as well. The global FRI theorem is
what shows that if the initial word is far from every degree-$<4$ codeword, then random
challenges and random query paths force either a detectable local inconsistency or a terminal
word that fails the final degree check.

\subsection{The complete proof obligations}

For this example the verifier's obligations are:
\begin{enumerate}[leftmargin=2em]
    \item accept $R_0$ only as a commitment to a length-$16$ word indexed by $D_0$;
    \item derive $\alpha_0$ after $R_0$ is fixed;
    \item accept $R_1$ only as a commitment to a length-$8$ word indexed by $D_1$;
    \item derive $\alpha_1$ after $R_1$ is fixed;
    \item receive the terminal constant $c_2=16$;
    \item derive query indices only after the complete commit transcript is fixed;
    \item verify every Merkle authentication path;
    \item check the round-$0$ folding equation;
    \item check the round-$1$ folding equation against $c_2$;
    \item check that the terminal description has degree less than $1$.
\end{enumerate}

The worked path verifies one randomly selected chain. A cryptographic instance repeats the
query phase at many independent indices and uses much larger fields and domains. The
arithmetic, however, is exactly the same as in the two equations checked above.
